%!TEX root = ../main.tex

\appendix
\chapter{Anhang}
\label{ch:anhang}

Der Anhang dokumentiert die Bestandteile der Untersuchung, die für die
Nachvollziehbarkeit der Ergebnisse erforderlich sind, den Lesefluss des
Haupttextes jedoch unterbrochen hätten: die wörtlichen Prompt-Vorlagen
(Abschnitt~\ref{apx:prompts}), die Aufgabentexte der neun \ac{ETL}-Aufgaben
(Abschnitt~\ref{apx:aufgaben}), den Aufgabentext der explizit gegliederten
Variante (Abschnitt~\ref{apx:komposition}) sowie den Verweis auf das
Repositorium mit Code, Daten und sämtlichen Ergebnisdatensätzen
(Abschnitt~\ref{apx:repo}).

\section{Prompt-Vorlagen}
\label{apx:prompts}

Die Vorlagen liegen als versionierte YAML-Dateien vor und bestehen jeweils aus
einer Systemanweisung und einer Nutzervorlage mit Platzhaltern, die zur Laufzeit
durch den Aufgabentext, die Eingabeverzeichnisse und den Ausgabepfad ersetzt
werden. Wiedergegeben ist der Wortlaut ohne die Platzhalterauflösung. Die
Vorlagen sind auf Englisch formuliert, da die Modelle in dieser Sprache
durchgehend trainiert und die Anbieter-Dokumentationen darauf ausgelegt sind.
Die Aufgabentexte selbst sind deutsch.

\subsection*{Code-Generierung}

\paragraph{v1 -- Zero-Shot:}
\emph{System:} \enquote{You are an experienced data engineer. Respond only with
executable Python code that uses pandas. No explanations, no markdown code
block, no text outside the code.}

\emph{Nutzervorlage:} \enquote{Write a self-contained Python script for the
following ETL task: \texttt{\{task\_description\}}. The input files are located
in: \texttt{\{input\_dir\}}. Write the result as a Parquet file to exactly:
\texttt{\{output\_path\}}. The script must run without any further input and
must produce the output file at the specified path.}

\paragraph{v2 -- Few-Shot:}
Systemanweisung wie v1. Die Nutzervorlage stellt der Aufgabe ein vollständig
gelöstes Beispiel voran, eingeleitet mit: \enquote{Here is one fully solved ETL
task as an example of the expected style. It is a different task from the one
you must solve; use it only to understand the format and quality of the expected
solution.} Das Beispiel behandelt eine Mitarbeitertabelle mit Leerzeichen in
Textspalten, fehlenden Abteilungen, als dekorierter Text abgelegten Gehältern
und den Wahrheitswerten \enquote{ja}/\enquote{nein}. Es umfasst Aufgabentext,
Eingabetabelle und die vollständige Pandas-Lösung. Die Beispielaufgabe ist
bewusst keine der neun evaluierten Aufgaben, damit sie keine Musterlösung
vorwegnimmt.

\paragraph{v3 -- Chain-of-Thought:}
\emph{System:} \enquote{You are an experienced data engineer. Respond only with
executable Python code (pandas). Record your reasoning as comments (with \#)
directly in the code. No text outside the code, no markdown.}

\emph{Nutzervorlage:} \enquote{Solve the following ETL task. Proceed step by
step: first break the task down into sub-steps (as comments), then implement
each step in code. \texttt{\{task\_description\}} \dots}

\subsection*{Direktverarbeitung}

Die drei Direkt-Vorlagen entsprechen den vorstehenden in Strategie und Beispiel,
verlangen jedoch die Ergebnistabelle statt des Codes.

\paragraph{d1 -- Zero-Shot:}
\emph{System:} \enquote{You are a precise data-processing engine. You receive
one or more input tables and a data transformation task. Apply the task directly
to the data and return ONLY the resulting table as CSV (comma-separated, with a
header row), wrapped in a single \texttt{csv} code block. Output every result
row. No explanations, no prose, no row index column.}

\emph{Nutzervorlage:} \enquote{Task: \texttt{\{task\_description\}}. Input data
(CSV): \texttt{\{data\}}. Return the COMPLETE resulting table as CSV inside one
\texttt{csv} code block. Include ALL result rows, no index column, no
commentary.}

\paragraph{d2 -- Few-Shot:}
Wie d1, ergänzt um dasselbe Beispiel wie v2, jedoch mit der Ergebnistabelle
statt des Lösungscodes als Vorbild.

\paragraph{d3 -- Chain-of-Thought:}
Wie d1, ergänzt um die Aufforderung: \enquote{Work through it step by step:
first identify the individual sub-steps the task consists of, then apply them to
the data one after another, so that each sub-step operates on the result of the
previous one.} Da die Antwort ausschließlich die Ergebnistabelle enthalten darf,
bleibt die Zerlegung anders als bei v3 modellintern.

\section{Aufgabentexte der neun ETL-Aufgaben}
\label{apx:aufgaben}

Wiedergegeben ist der Wortlaut, der als \texttt{task\_description} in die
Prompt-Vorlage eingesetzt wird. Für die schema-angereicherte Variante tritt
jeweils eine Beschreibung des Zielschemas hinzu, die hier aus Platzgründen
entfällt.

\subsection*{Datenbereinigung}

\begin{description}
    \item[Einfach (Leerzeichen und fehlende Werte):] Lies die Datei
    \texttt{customers\_raw.csv} ein. Entferne führende und nachfolgende
    Leerzeichen aus allen Textspalten. Ersetze fehlende Werte in der Spalte
    \texttt{country} durch \enquote{UNKNOWN}. Schreibe das Ergebnis als
    Parquet-Datei \texttt{cleaned\_customers.parquet}.

    \item[Mittel (Datumsformate):] Die Spalte \texttt{registered\_at} in
    \texttt{customers\_raw.csv} enthält Datumsangaben in unterschiedlichen
    Formaten (ISO-Datum, deutsches Format DD.MM.YYYY, US-Format Month DD YYYY
    und Unix-Timestamps). Normalisiere alle Werte auf das ISO-Format
    YYYY-MM-DD. Schreibe das Ergebnis als Parquet-Datei.

    \item[Schwer (semantische Vereinheitlichung):] Die Spalte \texttt{country}
    in \texttt{customers\_raw.csv} enthält Länder in verschiedenen Schreibweisen
    und Sprachen (z.\,B. \enquote{DE}, \enquote{Deutschland},
    \enquote{Germany}, \enquote{Deutshcland}). Vereinheitliche alle Werte auf
    den zweistelligen ISO-3166-1-alpha-2-Code (z.\,B. \enquote{DE}). Werte, die
    sich nicht eindeutig zuordnen lassen, sollen den Wert \enquote{UNKNOWN}
    erhalten. Schreibe das Ergebnis als Parquet-Datei.
\end{description}

\subsection*{Duplikaterkennung}

\begin{description}
    \item[Einfach (exakte Duplikate):] Lies \texttt{customers\_raw.csv} ein und
    entferne exakte Duplikate (Zeilen, in denen alle Spaltenwerte identisch
    sind). Behalte jeweils das erste Vorkommen. Schreibe das Ergebnis als
    Parquet-Datei.

    \item[Mittel (schlüsselbasiert):] Entferne in \texttt{customers\_raw.csv}
    Zeilen, in denen die \texttt{email} mehrfach vorkommt. Bei mehreren Zeilen
    mit gleicher E-Mail soll die Zeile mit dem jüngsten
    \texttt{registered\_at}-Wert behalten werden. Schreibe das Ergebnis als
    Parquet-Datei.

    \item[Schwer (unscharfe Duplikate):] In \texttt{customers\_raw.csv} befinden
    sich Fuzzy-Duplikate: gleiche Person, aber unterschiedliche Schreibweisen
    des Namens (Tippfehler, unterschiedliche Groß-/Kleinschreibung, mit/ohne
    Mittelinitial) und unterschiedliche E-Mail-Schreibweisen (Punkte vor dem @,
    unterschiedliche Groß-/Kleinschreibung). Erkenne diese Duplikate und
    entferne sie. Behalte jeweils die Zeile mit der vollständigsten Information.
    Schreibe das Ergebnis als Parquet-Datei.
\end{description}

\subsection*{Datentransformation}

\begin{description}
    \item[Einfach (Typkonvertierung):] Lies \texttt{products\_raw.csv} ein.
    Konvertiere die Spalte \texttt{price\_eur} von Text (z.\,B.
    \enquote{19,99}, \enquote{\texteuro 19.99}, \enquote{19.99 EUR}) in einen
    numerischen Wert (Float, Dezimaltrennzeichen Punkt). Konvertiere die Spalte
    \texttt{in\_stock} von Text (\enquote{ja}/\enquote{nein},
    \enquote{true}/\enquote{false}, \enquote{1}/\enquote{0}) in Boolean.
    Schreibe das Ergebnis als Parquet-Datei.

    \item[Mittel (abgeleitete Spalten):] Lies \texttt{orders\_raw.csv} ein.
    Berechne für jede Bestellzeile eine neue Spalte \texttt{total\_eur} als
    \texttt{quantity * unit\_price\_eur}. Extrahiere aus \texttt{ordered\_at}
    zusätzlich die Spalten \texttt{order\_year} (Integer) und
    \texttt{order\_month} (Integer, 1--12). Schreibe das Ergebnis als
    Parquet-Datei.

    \item[Schwer (Verknüpfung und Aggregation):] Verknüpfe
    \texttt{orders\_raw.csv}, \texttt{customers\_raw.csv} und
    \texttt{products\_raw.csv} über die jeweiligen ID-Spalten. Bereinige dabei
    die Daten so weit wie nötig (Typkonvertierungen, fehlende Werte sinnvoll
    behandeln). Aggregiere anschließend pro Land (\texttt{country\_code},
    ISO-2) und Produktkategorie (\texttt{category}) den Gesamtumsatz
    (\texttt{total\_revenue\_eur = sum(quantity * unit\_price\_eur)}) und die
    Anzahl Bestellungen (\texttt{order\_count}). Sortiere absteigend nach
    Gesamtumsatz. Schreibe das Ergebnis als Parquet-Datei.
\end{description}

\section{Aufgabentext der explizit gegliederten Variante}
\label{apx:komposition}

Dieser Text stellt dieselbe Verarbeitung wie die schwere Transformationsaufgabe,
benennt die erforderlichen Teilschritte jedoch ausdrücklich. Abschnitt
\ref{subsec:komposition} wertet die damit erzielten Ergebnisse aus.

\begin{quote}
    Erzeuge aus den drei Rohtabellen \texttt{customers\_raw.csv},
    \texttt{products\_raw.csv} und \texttt{orders\_raw.csv} eine
    Umsatzauswertung. Führe dazu die folgenden Schritte in dieser Reihenfolge
    aus:

    \begin{enumerate}
        \item Entferne in \texttt{customers\_raw.csv} führende und nachfolgende
        Leerzeichen aus allen Textspalten.
        \item Vereinheitliche die Spalte \texttt{country} in
        \texttt{customers\_raw.csv} auf den zweistelligen
        ISO-3166-1-alpha-2-Code und lege sie als \texttt{country\_code} ab. Die
        Rohwerte liegen in unterschiedlichen Sprachen, als Abkürzungen und mit
        Tippfehlern vor (z.\,B. \enquote{Deutschland}, \enquote{Germany},
        \enquote{GER}, \enquote{Deutshcland} $\rightarrow$ \enquote{DE}). Werte,
        die sich nicht eindeutig zuordnen lassen, sowie fehlende Werte erhalten
        \enquote{UNKNOWN}.
        \item Entferne aus \texttt{customers\_raw.csv} Datensätze mit mehrfach
        vorkommender \texttt{customer\_id} und behalte jeweils das erste
        Vorkommen.
        \item Konvertiere in \texttt{products\_raw.csv} die Spalte
        \texttt{price\_eur} (Text wie \enquote{19,99}, \enquote{19.99 EUR}) in
        einen Float und \texttt{in\_stock} in einen Boolean.
        \item Verknüpfe \texttt{orders\_raw.csv} mit
        \texttt{customers\_raw.csv} über \texttt{customer\_id} und mit
        \texttt{products\_raw.csv} über \texttt{product\_id}. Alle Bestellzeilen
        aus \texttt{orders\_raw.csv} bleiben dabei erhalten.
        \item Aggregiere pro Land (\texttt{country\_code}) und Produktkategorie
        (\texttt{category}) den Gesamtumsatz \texttt{total\_revenue\_eur =
        sum(quantity * unit\_price\_eur)}, auf zwei Nachkommastellen gerundet,
        sowie die Anzahl Bestellungen \texttt{order\_count}. Die Spalten
        \texttt{quantity} und \texttt{unit\_price\_eur} stammen beide aus
        \texttt{orders\_raw.csv}; der Preis aus \texttt{products\_raw.csv} wird
        für die Umsatzberechnung nicht verwendet. Sortiere absteigend nach
        \texttt{total\_revenue\_eur}.
    \end{enumerate}
\end{quote}

\section{Repositorium}
\label{apx:repo}

Sämtlicher Code, die Erzeugungsvorschrift des Datensatzes, die regelbasierte
Vergleichspipeline, die Prompt-Vorlagen sowie alle Ergebnisdatensätze der 711
Modellaufrufe stehen unter
\url{https://github.com/Lu1s198/Bachelorarbeit} zur Verfügung. Die Struktur
folgt den Kapiteln dieser Arbeit:

\begin{description}
    \item[\texttt{project/src/dataset/}] Erzeugung des synthetischen Datensatzes
    einschließlich der Aufgabendefinitionen (Kapitel~\ref{ch:datensatz}).
    \item[\texttt{project/src/baseline/}] Regelbasierte Vergleichspipeline
    (Abschnitt~\ref{sec:baseline}).
    \item[\texttt{project/src/llm/} und \texttt{project/prompts/}] Anbindung der
    Modelle und die Prompt-Vorlagen (Abschnitt~\ref{sec:llm-integration}).
    \item[\texttt{project/src/experiments/}] Die Ausführungsschichten der
    einzelnen Untersuchungsebenen (Kapitel~\ref{ch:durchfuehrung}).
    \item[\texttt{project/src/evaluation/}] Umsetzung der Evaluationskriterien
    aus Abschnitt~\ref{sec:evaluationskriterien}.
    \item[\texttt{project/results/}] Die Ergebnisdatensätze je Lauf sowie die
    erzeugten Skripte und Ergebnistabellen.
\end{description}

Die berichteten Ergebnisse stammen sämtlich aus dem Datensatz mit dem Startwert
2 beziehungsweise dessen halbierter Fassung mit dem Startwert 102.
